% 第8章 デリバティブ(資料8)
\chapter{デリバティブ}

本章から講義の後半に入り，\textbf{デリバティブ}(金融派生商品，derivatives)を
扱う。本章ではまず，先渡・先物，スワップ，オプションという代表的なデリバティブの
仕組みを学び，ペイオフと正味損益という評価の視点，およびデリバティブを用いた
取引戦略を整理する。価格付けの理論は次章以降で扱う。

\section{デリバティブ(金融派生商品)}

\begin{definition}[デリバティブ]
\textbf{デリバティブ(金融派生商品)}とは，伝統的な金融商品(株式，債券，金利，
為替など)から派生して生まれた金融商品のことである。\textbf{派生証券}とは，
既存の金融商品の価格・収益率，キャッシュフロー等に基づいて将来のキャッシュ
フローが決まる証券である。
\end{definition}

デリバティブから派生するデリバティブ，というものもある。

\begin{example}[江戸時代の米市場]
江戸時代，大阪堂島の米市場では，米の収穫前に取引価格を決める取引が行われて
いた。これは米の現物取引から派生して生まれた取引で，デリバティブのはしりと
考えられている。
\end{example}

デリバティブの名称は「対象+デリバティブの種類」で表される。
\begin{itemize}
\item \textbf{先物}：株価指数先物，債券先物，金利先物，…
\item \textbf{オプション}：株式オプション，債券オプション，金利オプション，
通貨オプション，…
\item \textbf{スワップ}：金利スワップ，通貨スワップ，エクイティスワップ，…
\end{itemize}

取引形態としては取引所取引と相対取引(店頭取引)がある。取引所で取引される
デリバティブは標準化されている。一方，当事者同士で相対取引されるデリバティブは，
対象も内容も多様である。

\section{先渡契約と先物契約}

先渡と先物は似ているが，微妙に異なる。

\subsection{先渡契約}

\begin{example}
ある投資家が，(今は元手がないが1年後に100万円入るので)1年後にA社の株式を
1株購入したい。今の株価は90万円である。1年後の株価が100万円を越えると購入
できない。そこで，現時点で「1年後に95万円でA社の株式を購入する」という契約を
結んでおけば，投資家は1年後に確実に株式を入手できる。しかし，1年後の価格が
85万円になったとしても，95万円で購入しなければならない\emph{義務}が生じる。
\end{example}

こういった契約を\textbf{先渡(さきわたし，フォワード)契約}という。購入側を
\textbf{ロングポジション}，売却側を\textbf{ショートポジション}という。

\begin{definition}[先渡契約]
ある資産(\textbf{原資産})を，将来の決められた時点(\textbf{満期})に，現時点で
決められた価格(\textbf{先渡価格})で売買することを定めた契約。
\end{definition}

上の例では，原資産はA社の株式，満期は1年後，先渡価格は95万円である。1年後に
実際にA社の株価がいくらになるかはわからないので，
\begin{itemize}
\item 1年後の株価が95万円より高ければ，買う側が儲かる。
\item 1年後の株価が95万円より低ければ，売る側が儲かる。
\end{itemize}
すなわち，先渡契約では儲かることもあるし，損することもある。

\begin{keypoint}[先渡価格の決まり方]
先渡では，\textbf{契約時点(現時点)で金銭のやり取りはない}(今は契約を結ぶ
だけ)。ということは，\textbf{この契約の価値は現時点でゼロ}である。つまり，
\emph{現在価値がゼロになるように先渡価格が決まる}。この考え方は金融取引に
共通である。そして，1年後には\emph{必ず}株式と代金を交換しなければならない。
\end{keypoint}

\subsection{先物契約}

\begin{definition}[先物契約]
\textbf{先物(フューチャーズ)契約}とは，概念上は，取引所で取引される定型化
された先渡契約である。ただし正確に言うと，先渡契約とはキャッシュフローが
異なる。すなわち，(満期で)現物の受渡しをせず，\textbf{差金決済}を行う
(先物取引の決済で使われる価格を\textbf{先物価格}という)。
\end{definition}

\begin{itemize}
\item \textbf{差金決済}：現物とキャッシュを交換するのが先渡契約の決済方法で
あるのに対し，購入価格と決済時点の先物価格の差額をやり取りする決済方法。
\item 先物契約では，毎日\textbf{値洗い}し，差金決済を行う。値洗いとは，適宜
そのときの時価(市場価格)で評価していくことである。この場合の差金は，
(当日の先物価格)$-$(前日の先物価格)である。
\item 先物価格は，\textbf{満期日に現物(原資産)価格と一致}する。
\end{itemize}

\begin{example}[先渡契約と先物契約のキャッシュフロー]
6/15が契約日，6/19が先物の満期日とし，原資産価格と先物価格が次のように推移した
とする。
\vspace{0.9em}
\begin{center}
\small
\begin{tabular}{llcc|ccc|cc}
\toprule
 & & 原資産 & 先物 & \multicolumn{3}{c|}{先渡契約} & \multicolumn{2}{c}{先物契約}\\
 & 日付 & 価格 & 価格 & 支払額 & 受取資産の価格 & 計 & 差金決済 & 累積額\\
\midrule
契約日 & 6/15 & 100 & 102 & & & & &\\
 & 6/16 & 98 & 100 & & & & $-2$ & $-2$\\
 & 6/17 & 101 & 102 & & & & $2$ & $0$\\
 & 6/18 & 104 & 105 & & & & $3$ & $3$\\
満期 & 6/19 & 101 & 101 & 102 & 101 & $-1$ & $-4$ & $-1$\\
\bottomrule
\end{tabular}
\end{center}
\vspace{0.9em}
先渡契約(先渡価格は契約日の先物価格102円)では，満期に102円を支払って時価
101円の資産を受け取るので，決済の価値は $-1$ 円である。一方，先物契約では毎日
差金決済を行い，満期日に締める。累積額は $-2+2+3-4 = -1$ 円で，最終的な累積額は
先渡と同じになるが，キャッシュフローの発生タイミングが異なる。このケースでは
(ほぼ)裁定取引ができてしまう(詳細は略)。
\end{example}

満期以前に先物を清算するには，\textbf{反対売買}を行う。
\begin{itemize}
\item ロングポジションの清算では，保有量と同量のショートポジションを追加する
ことで，ロングポジションを解消する。その後のロングポジションによるキャッシュの
交換は，ショートポジションによるキャッシュの交換と相殺される，と考えればよい。
\item ショートポジションの清算では，保有量と同量のロングポジションを追加する
ことで，ショートポジションを解消する。
\end{itemize}

先物契約はいろいろな商品にあり，リスクヘッジに活用される。
\begin{itemize}
\item \textbf{金融商品先物}：株価指数先物，国債先物，金利先物，…
\item \textbf{商品先物}：金先物，白金先物，原油先物，とうもろこし先物，…
\end{itemize}

\section{スワップ}

\begin{definition}[スワップ契約]
\textbf{スワップ契約}とは，将来のキャッシュフロー列を別のキャッシュフロー列と
交換する契約である。
\end{definition}

名称は「対象+スワップ」で表示される。
\begin{itemize}
\item \textbf{金利スワップ}：異なる金利による利払の交換。固定金利と変動金利の
キャッシュフロー列の交換が最も多い。
\item \textbf{通貨スワップ}：異なる通貨のキャッシュフロー列を交換。
\item \textbf{ベーシススワップ}：異なる変動金利によるキャッシュフロー列を交換。
\end{itemize}

スワップ契約は店頭取引である。定型化された契約が多いが，オーダーメイドにより
キャッシュフローを自由に設計できる。

\subsection{金利スワップ}

最も典型的なのが金利スワップであり，定期的な\textbf{固定払}と\textbf{変動払}を
交換する契約が主流である。代表例では，同じ元本に対する固定金利払と変動金利払の
キャッシュフロー列を交換する(例えば「固定金利受け・変動金利払い」)。

\paragraph{変動金利払のキャッシュフローの意味。}
変動払のキャッシュフロー列は，\textbf{短期借換}のキャッシュフローとして理解
できる。一期間の貸付は「元本の貸付，元本+利息の受取」であり，短期借換は
「元本の貸付，元本+利息の受取」の繰り返しである。このとき，同時点の元本の
貸付と返済を相殺すると，実質的に残るキャッシュフローは「各利払日の利息+最終時点の
元本」，すなわち\textbf{変動利付債のキャッシュフロー}になる。

\paragraph{金利スワップのキャッシュフローの合成。}
「変動金利のキャッシュフロー(変動利付債)」の受けと「固定金利のキャッシュ
フロー(固定利付債)」の払いを組み合わせると，元本が同額ならば最初と最後の
元本キャッシュフローが相殺され，残ったキャッシュフロー列は金利スワップのものと
同じになる。つまり，金利スワップは「変動利付債ロング+固定利付債ショート」
(またはその逆)と分解できる。

金利スワップに関する用語と性質を整理する。
\begin{itemize}
\item キャッシュフローを決めるには元本が必要である。これを\textbf{想定元本}
(nominal principal) と呼ぶ。交換するのは金利部分のみで，元本は交換しない
(だから“想定”元本)。
\item 日本の債券は固定払が多いので，債券を購入して固定金利を受け取り，金利
スワップ(固定金利払・変動金利受，payers swap)を契約すると，金利リスクは低下
する。
\item 金融機関の資金調達は変動金利と考えてよいので，基本は変動金利払である。
一方，固定利付債を保有すると固定金利受となる。上記スワップを締結すると変動金利受
に変わるので，変動金利払とマッチする。\textbf{支払と受取のタイプを合わせた方が，
リスクが低下する}。
\end{itemize}

\begin{example}[金利スワップの例 (1)：債券投資における金利リスクヘッジ]
A社は，B社の社債(元々満期5年，クーポンレート3\%，残存2年)を1億円で購入し，
半年ごとにクーポン $1\text{億円}\times 3\%/2$ を受け取る。A社の資金調達は
変動金利(LIBOR)による(半年ごとの借換，2年間。市場に半年ごとに
$1\text{億円}\times(1+\text{LIBOR}/2)$ を支払い，1億円を借り換える)。そこで
A社は，C社と金利スワップ(満期2年)を契約し，固定払(半年ごと，
$1\text{億円}\times 3\%/2$)・変動受(半年ごと，
$1\text{億円}\times\text{LIBOR}/2$)とする。これにより，B社債からの固定受と
スワップの固定払が相殺され，スワップの変動受が資金調達の変動払とマッチして，
金利リスクがヘッジされる。

なお，スワップ取引の固定金利のことを\textbf{スワップ金利(スワップレート)}と
いう。(実際にはLIBORと固定金利3\%には乖離がありうる。)
\end{example}

\begin{example}[金利スワップの例 (2)：将来の金利上昇への備え]
A社はB社から変動金利で10億円を5年間借入れている(変動払：半年ごと
$10\text{億円}\times\text{LIBOR}/2$)。現在，6ヶ月LIBOR $=2\%$，5年スワップ金利
$=2.5\%$ とする。将来の金利上昇による資金調達コスト増のリスクを回避するため，
A社はC社と満期5年の金利スワップを契約し，固定払(半年ごと
$10\text{億円}\times 2.5\%/2$)・変動受(半年ごと
$10\text{億円}\times\text{LIBOR}/2$)とする。これで，実質的な支払は固定金利
2.5\%に固定され，金利上昇リスクは回避される。ただし，将来のLIBORが2.5\%を
下回ると，(固定してしまった分)損になる。
\end{example}

\vspace{18pt}

\subsection{通貨スワップ}

\begin{definition}[通貨スワップ]
\textbf{通貨スワップ}とは，通貨の異なるキャッシュフローを交換するスワップである。
例：円3ヶ月ON金利とドル固定金利の交換。元本の交換をすることが多い。元本を交換
しない場合，\textbf{クーポン・スワップ}とも呼ぶ。
\end{definition}

ここで3ヶ月ON (OverNight) 金利とは，3ヶ月間のON金利による運用成果と同等の
変動金利のことである。

\begin{example}[通貨スワップの例]
ドルの調達を実質的に円の調達に変換する($1\$ \approx 160$円として)。企業が
債権者から1億ドルを2年間，金利2\%で調達しているとする。銀行と，元本1億ドル，
金利2\%(年利)，スワップレート160円/ドル，満期2年，半年払の通貨スワップを
契約すると，キャッシュフローは次のようになる。
\begin{itemize}
\item $t=0$：企業は債権者から1億ドルを受け取り，それを銀行に渡して160億円を
受け取る(元本交換)。
\item $t=0.5,1,1.5,2$年：企業は銀行に1.6億円(160億円$\times$2\%/2)を支払い，
銀行から100万ドル(1億ドル$\times$2\%/2)を受け取って債権者への利払に充てる。
\item $t=2$年：企業は銀行に160億円を返し，1億ドルを受け取って債権者に償還する
(元本再交換)。
\end{itemize}
これにより，企業のキャッシュフローは実質的に「160億円を金利2\%で調達」した形に
なる。
\end{example}

\vspace{18pt}

\subsection{ベーシス・スワップなど}

\textbf{ベーシス・スワップ}とは，異なる変動金利のキャッシュフローの交換である。
例1：円3ヶ月ON金利と円3ヶ月TIBORの交換。例2：円3ヶ月ON金利とドル3ヶ月ON金利の
交換。通貨が異なる場合は，元本を交換することが多い。

あるスワップは，他のスワップの組み合わせでも表現可能である。例えば，
\begin{align*}
&\text{元本交換を伴う通貨スワップ(円3ヶ月ON金利 vs ドル固定金利)}\\
&= \text{元本交換を伴うベーシス・スワップ(円3ヶ月ON金利 vs ドル3ヶ月ON金利)}\\
&\quad + \text{金利スワップ(ドル3ヶ月ON金利 vs ドル固定金利)}
\end{align*}
である。

\subsection{クレジット・デフォルト・スワップ(CDS)}

\textbf{CDS} (Credit Default Swap) は，参照資産Cの保有者にとって保険のような
ものである(ここでは single name CDS を説明する)。

\begin{itemize}
\item 契約の要素：参照組織・参照債務，想定元本，満期(予定終了日)，固定金利
(スプレッド)。
\item プロテクションの買い手は，固定金利(=想定元本$\times$スプレッド)を
典型的には四半期毎・後払いで支払う(\textbf{プレミアム・レッグ})。
\item \textbf{クレジット・イベント(CE)}(バンクラプシー，支払不履行，
リストラクチャリング)が発生すると，売り手は損失を補償する(参照債務の引渡し等，
\textbf{デフォルト・レッグ})。
\end{itemize}

購入者側のキャッシュフローは次のとおりである。
\begin{itemize}
\item 満期までデフォルトが発生しないとき：開始から満期まで，プレミアム・レッグ
(スプレッドの支払い)のみが発生する。
\item 満期までにデフォルトが発生するとき：デフォルト発生時点までプレミアム・
レッグを支払い，発生時点でデフォルト・レッグ(補償)を受け取って契約は終了する。
\end{itemize}

\section{オプション}

\subsection{オプションの例}

\begin{example}[例1：買う権利]
「A社の株式(現在56万円)を，1年後に57万円で\textbf{買うことができる権利}」を
考える。
\begin{itemize}
\item 1年後に株価が60万円になった場合：権利行使すれば57万円で買えるので，市場で
買うよりも3万円の得。
\item 1年後に株価が57万円になった場合：権利行使しても現物を買っても57万円で
買えるので，どちらでも同じ。
\item 1年後に株価が50万円になった場合：権利行使せず，市場で50万円で買えば，
損も得もしない。
\end{itemize}
満期におけるこの権利(オプション)の価値は，満期の株価を $S_T$，行使価格を
$K=57$万円として
\[
V_T = \max(S_T - K,\ 0)
\]
である($S_T$ が $K$ を超えた分だけ価値をもつ折れ線グラフになる)。
\end{example}

\begin{example}[例2：売る権利]
「A社の株式(現在56万円)を，1年後に57万円で\textbf{売ることができる権利}」を
考える。
\begin{itemize}
\item 1年後に株価が50万円になった場合：権利行使すれば57万円で売れるので，市場で
売るよりも7万円の得。
\item 1年後に株価が57万円になった場合：どちらでも同じ。
\item 1年後に株価が60万円になった場合：権利行使せず，市場で60万円で売れば，
損も得もしない。
\end{itemize}
満期におけるこの権利の価値は
\[
V_T = \max(K - S_T,\ 0)
\]
である。
\end{example}

\vspace{18pt}

\subsection{オプションの定義}

\begin{definition}[オプション]
\textbf{オプション}とは，「将来に，ある証券をある価格で取引する権利」を表章する
証券である。
\end{definition}

\begin{itemize}
\item オプションを購入して得られるものは「\textbf{権利}」であり，購入した側には
しなければならない「義務」は生じない(使って得になる場合だけ使えばよい)。
\item 一方で，オプションを売却した側には，オプション購入者の要求に応じる
「\textbf{義務}」が生じる。
\item 買い手は\textbf{ロングポジション}，売り手は\textbf{ショートポジション}。
\item この権利を得るために，買い手は代金(\textbf{オプションプレミアム})を
売り手に支払う。この将来の権利には価値があるので，契約時に現在価値を支払う
(金融取引の原則に沿っている)。
\item 契約時に現在価値の分だけ代金を支払うので，権利を行使しなければ，支払った
代金分だけ損をする。
\end{itemize}

オプションのキーワードは，\textbf{原資産}，\textbf{権利の種類}(買う権利か，
売る権利か)，\textbf{権利行使日}，\textbf{権利行使価格}である。

\subsection{コールとプット，権利行使日}

\begin{definition}[コール・オプションとプット・オプション]
\begin{itemize}
\item \textbf{コール・オプション}：ある資産(原資産)を，ある特定の日(権利
行使日)に，ある価格(権利行使価格)で\textbf{買う権利}。
\item \textbf{プット・オプション}：ある資産(原資産)を，ある特定の日(権利
行使日)に，ある価格(権利行使価格)で\textbf{売る権利}。
\end{itemize}
行使すると損になる時は，行使しなくても良い。
\end{definition}

権利行使日による分類は次のとおりである。
\begin{itemize}
\item \textbf{ヨーロピアン・オプション}：権利行使はオプションの満期日のみ可能。
\item \textbf{アメリカン・オプション}：権利行使はオプションの満期日まで，
いつでも可能。
\end{itemize}
名称はこれらを組み合わせる(例：ヨーロピアン・コール・オプション，ヨーロピアン・
プット・オプション，アメリカン・プット・オプション，…)。なお，アジアン・
オプション(平均値オプション)，ロシアン・オプション(最大値オプション)なども
ある。

\begin{example}[行使時の価値]
①3ヶ月後に日経平均を65,000円で売る権利(プット・オプション)：行使時の受取額は
株価が65,000円を下回った分であり，オプション価値は $\max(65{,}000-S_T, 0)$。
②1ヶ月後にA社株を500円で買う権利(コール・オプション)：行使時の実質的な
支払額は $\min(S_T, 500)$ 円で済み，オプション価値は $\max(S_T-500, 0)$。
\end{example}

\vspace{18pt}

\subsection{ITM，OTM，ATM}

行使価格と原資産価格の相対的な位置を表す用語がある。

\begin{itemize}
\item \textbf{In The Money (ITM，イン・ザ・マネー)}：もし仮に現時点で
オプションを行使したら，利益が出る状態。
\item \textbf{Out of The Money (OTM，アウト・オブ・ザ・マネー)}：もし仮に
現時点でオプションを行使したら，利益が出ない状態。
\item \textbf{At The Money (ATM，アット・ザ・マネー)}：原資産の現在価値が
行使価格に等しい状態。$S(t) = K$($K$：行使価格)。
\item \textbf{Near The Money(ニア・ザ・マネー)}：ATM付近の状態。
\end{itemize}

コールオプションならば，ITMは $S(t) > K$，OTMは $S(t) < K$。プットオプション
ならば，ITMは $S(t) < K$，OTMは $S(t) > K$ である。

\section{ペイオフと正味損益}

\subsection{ペイオフ関数}

満期 $T$ における原資産価格 $S(T)$ の関数として，オプションの満期時の受取
(\textbf{ペイオフ})を描いたものをペイオフ関数という。これらのキャッシュ
フローは，$S(T)$ に依存する確率変数である。

\begin{itemize}
\item コール・ロング：$\max(S(T)-K, 0)$(行使価格 $K$ から右上がりの折れ線)。
\item コール・ショート：$-\max(S(T)-K, 0)$(行使価格 $K$ から右下がり)。
\item プット・ロング：$\max(K-S(T), 0)$($S(T)=0$ で最大値 $K$，$K$ でゼロに
なる折れ線)。
\item プット・ショート：$-\max(K-S(T), 0)$($S(T)=0$ で最小値 $-K$)。
\end{itemize}

\begin{figure}[htbp]
\centering
\begin{tikzpicture}[scale=0.62]
% call long
\begin{scope}
\draw[-{Stealth}] (0,0) -- (5,0) node[below]{$S(T)$};
\draw[-{Stealth}] (0,-1.6) -- (0,2.4) node[above]{ペイオフ};
\draw[very thick, blue] (0,0) -- (2.5,0) -- (4.6,2.1);
\node[below] at (2.5,0) {$K$};
\node at (2.3,2.0) {\small コール・ロング};
\end{scope}
% call short
\begin{scope}[xshift=6.8cm]
\draw[-{Stealth}] (0,0) -- (5,0) node[below]{$S(T)$};
\draw[-{Stealth}] (0,-2.4) -- (0,1.6) node[above]{ペイオフ};
\draw[very thick, red] (0,0) -- (2.5,0) -- (4.6,-2.1);
\node[above] at (2.5,0) {$K$};
\node at (2.6,1.1) {\small コール・ショート};
\end{scope}
% put long
\begin{scope}[yshift=-6.2cm]
\draw[-{Stealth}] (0,0) -- (5,0) node[below]{$S(T)$};
\draw[-{Stealth}] (0,-1.6) -- (0,2.4) node[above]{ペイオフ};
\draw[very thick, blue] (0,1.8) -- (2.5,0) -- (4.6,0);
\node[below] at (2.5,0) {$K$};
\node[left] at (0,1.8) {$K$};
\node at (3.3,1.2) {\small プット・ロング};
\end{scope}
% put short
\begin{scope}[xshift=6.8cm, yshift=-6.2cm]
\draw[-{Stealth}] (0,0) -- (5,0) node[below]{$S(T)$};
\draw[-{Stealth}] (0,-2.4) -- (0,1.6) node[above]{ペイオフ};
\draw[very thick, red] (0,-1.8) -- (2.5,0) -- (4.6,0);
\node[above] at (2.5,0) {$K$};
\node[left] at (0,-1.8) {$-K$};
\node at (3.4,1.0) {\small プット・ショート};
\end{scope}
\end{tikzpicture}
\caption{コール/プットオプション(ロング，ショート)のペイオフ関数。}
\label{fig:payoff}
\end{figure}

\subsection{正味損益}

満期における\textbf{正味損益}は，ペイオフに当初支払った(受け取った)
オプション・プレミアムの満期での価値を加味したものである。
\begin{itemize}
\item コール・ロングの正味損益は，ペイオフ関数を「$C(0)$(プレミアム)の満期
での価値」だけ下にシフトしたもの。
\item コール・ショートの正味損益は，ペイオフ関数を同じ分だけ上にシフトしたもの。
\item プットについても同様である。
\end{itemize}

\begin{keypoint}[正味損益の重要性]
デリバティブの効用は，正味損益で考えるのが妥当である。ペイオフ関数はもちろん
重要だが，そのペイオフ関数を得るためにコスト(プレミアム払い)がかかっている。
なので，「コスト込みの評価 $=$ 正味損益による評価」が重要である。

では，肝心のコスト($=$プレミアム)はどのように算出するか?それが
\textbf{デリバティブの価格付け}であり，次章以降で説明する。
\end{keypoint}

\section{デリバティブを用いた取引戦略}

\subsection{ヘッジ}

\begin{definition}[ヘッジ]
\textbf{ヘッジ}とは，オプションとその原資産を組み合わせたポジションのことである。
\begin{itemize}
\item \textbf{1対1ヘッジ}：原資産1単位ロング+コール1単位ショート。
\item \textbf{1対2ヘッジ}：原資産1単位ロング+コール2単位ショート。
\item \textbf{1対1逆ヘッジ}：原資産1単位ショート+コール1単位ロング。
\end{itemize}
\end{definition}

満期の正味損益は次のようになる。
\begin{itemize}
\item \textbf{1対1ヘッジ}:$S(T)<K$ では原資産ロングの損益(右上がり)，
$S(T)>K$ ではコールショートの損失と相殺して水平になる。この形は
\emph{プットショートと同じ形}だが，水準は一般に異なる(後の講義でこれに関係する
話(プット・コール・パリティ)が出てくる)。
\item \textbf{1対2ヘッジ}:$S(T)<K$ では右上がり，$S(T)>K$ ではコール2単位の
ショートが効いて右下がりになる(山型)。
\item \textbf{1対1逆ヘッジ}：1対1ヘッジの逆で，\emph{プットロングと同じ形}だが，
水準は一般に異なる。
\end{itemize}

\subsection{スプレッド}

\begin{definition}[スプレッド]
\textbf{スプレッド}とは，行使価格あるいは満期が異なる(同じ原資産の)コール
オプション(またはプットオプション)同士を組み合わせたポジションのことである。
\begin{itemize}
\item \textbf{バーティカル・スプレッド}：行使価格が異なるオプション同単位を
組み合わせたポジション。
  \begin{itemize}
  \item 強気(ブル)：低い行使価格 $L$ のコールロングと高い行使価格 $K$ の
  コールショート。\textbf{ブル・スプレッド戦略}という。
  \item 弱気(ベア)：低い行使価格 $L$ のコールショートと高い行使価格 $K$ の
  コールロング。\textbf{ベア・スプレッド戦略}という。
  \item バタフライ：(低い)強気と(高い)弱気の組み合わせ。
  \end{itemize}
\item \textbf{ホリゾンタル・スプレッド}：満期が異なるオプションを組み合わせた
ポジション。
\end{itemize}
\end{definition}

満期の正味損益は次のようになる。
\begin{itemize}
\item \textbf{バーティカル・ブル・コール・スプレッド}:$S(T)<L$ で一定の小さな
損失，$L<S(T)<K$ で右上がり，$S(T)>K$ で一定の利益。相場上昇を見込む戦略。
\item \textbf{バーティカル・ベア・コール・スプレッド}：上の逆で，$S(T)<L$ で
一定の利益，$S(T)>K$ で一定の損失。相場下落を見込む戦略。
\item \textbf{ロング・バタフライ・スプレッド}：行使価格 $L$ のコールロング，
行使価格 $M$ のコールショート2単位，行使価格 $K$ のコールロング($L<M<K$)から
なり，$S(T)=M$ 付近で利益が最大になる山型。相場が動かないことを見込む戦略。
\item \textbf{ショート・バタフライ・スプレッド}：ロング・バタフライの上下を
ひっくり返したもの。個々のデリバティブのロングをショートに，ショートをロングに
すれば構築できる。
\end{itemize}

\subsection{コンビネーション}

\begin{definition}[コンビネーション]
\textbf{コンビネーション}とは，同じ原資産のコールオプションとプットオプションを
組み合わせたポジションのことである。

\textbf{ストラドル}：同じ原資産で行使価格が同じコールオプションとプット
オプション同単位を組み合わせたポジション。
\end{definition}

\textbf{ボトムストラドル}(コールロング+プットロング)の正味損益はV字型で，
$S(T)=K$ 付近で最大の損失(支払ったプレミアム分)，そこから離れるほど利益が
増える。相場が大きく動くことを見込む戦略である。

\subsection{レバレッジ効果}

デリバティブでは，(想定)元本と同額のキャッシュを取引することはあまりない。
一方，コストなしに自由に取引できるわけではなく，(想定)元本額に比べて少額の
\textbf{証拠金}($\sim$担保)を積んでおく必要がある。

\begin{keypoint}[レバレッジ効果]
少額の証拠金を積むことで，高額の元本による取引と同等の経済効果を得ることが
できる。これは\textbf{レバレッジ効果}と呼ばれる。このため，デリバティブは投機
目的で投資されることもある。
\end{keypoint}

% ---------------- 演習問題 ----------------
\section{演習問題}

\begin{exercise}[問題1]
本日の日経平均は60,000円である。投資家Aは，今後3ヶ月間，日経平均は今の価格から
$\pm 2{,}000$円のレンジに収まると読んでいる。そこで，1ヶ月後にこのレンジで一定額の
収益を得られる戦略を，ヨーロピアン・コール・オプションとヨーロピアン・プット・
オプション1単位ずつを組み合わせて作ろうと考えた。満期1ヶ月のどのような
オプションを組み合わせればよいか。また，そのときの満期時の正味損益関数を示し
なさい。ただし，このレンジを離れるにつれて収益が低減(場合によっては損失が
発生)するのは仕方がないものとする。

※ 正味損益関数は，満期のペイオフ関数に，当初やり取りするオプション・プレミアムの
効果も考慮した価値。プレミアムは計5,400円程度。
\end{exercise}

\begin{answerbox}
\textbf{考え方。}「レンジ内で一定額の収益」を得たいのだから，レンジ内でペイオフが
一定(フラット)になる組み合わせを探す。オプションの\emph{売り}はプレミアム収入を
生むので，レンジ内でどのオプションも行使されない(ペイオフがゼロの)ポジションを
売りで組めば，プレミアム分が確定収益になる。

\medskip
\textbf{解答。}
「\textbf{行使価格58,000円のプットオプション1単位の売り}」と，
「\textbf{行使価格62,000円のコールオプション1単位の売り}」を組み合わせれば
よい。

\begin{itemize}
\item 満期の日経平均が $58{,}000 \leq S_T \leq 62{,}000$ のレンジに収まれば，
プットもコールも行使されず，ペイオフはゼロである。
\item プットとコールの売りによるペイオフ関数だけでは，満期時にプラスにならない。
しかし，このポジションを構築する際にプットとコールを\emph{売却}しているので，
それらのオプション・プレミアム(計5,400円程度)を既に得ている。
\item それらの満期時点での価値(こういうときは無リスク資産で運用すると考える)を
加えることで，正味損益関数はレンジ内で一定のプラス(約5,400円強)になる。
\end{itemize}

正味損益関数の形は，$S_T < 58{,}000$ では傾き $+1$ で下がっていき(プットが
行使される)，$58{,}000\leq S_T\leq 62{,}000$ では一定のプラス，
$S_T > 62{,}000$ では傾き $-1$ で下がっていく(コールが行使される)，
台形の上辺のような形である。レンジを大きく外れると損失が発生する。

\medskip
\textbf{補足：ストラングル。}このような戦略は，\textbf{ストラングル}と呼ばれる
コンビネーション戦略の一種である。ストラングルとは，同じ原資産で\emph{行使価格が
異なる}コールオプションとプットオプション同単位を組み合わせたポジションであり，
本問のものは\textbf{ショート・ストラングル戦略}である。上下を逆転させた戦略も
あり，\textbf{ロング・ストラングル戦略}と呼ばれる。
\end{answerbox}

\begin{exercise}[問題2]
問題1の投資家Aはその後再考し，日経平均がもしも大幅に上昇しても巨額の損失を
被ることは避けたいと考えた(日経平均が大幅に下落することはないと固く信じて
いる)。その場合，どのようなポジションを追加すればよいか。満期3ヶ月の1種類の
コール・オプションの契約(売り or 買い)1単位を問題1の戦略に追加することで，
望ましいポジションを作りなさい。また，そのときの満期時の正味損益関数の形を示し
なさい。

※ このとき，正味損益関数のピーク(最大値)はどうなるか考えてみること。
\end{exercise}

\begin{answerbox}
\textbf{考え方。}問題1のポジションの弱点は，$S_T$ が大きく上昇したときに
コールショートの損失が\emph{無限に}拡大することである。上昇側の損失を頭打ちに
するには，より高い行使価格のコールを\emph{買って}おけばよい(ブル側をスプレッドに
変える発想)。

\medskip
\textbf{解答。}
問題1のポジション(行使価格58,000円のプット売り+行使価格62,000円のコール
売り)に，「\textbf{行使価格がより高いコールオプション1単位の買い}」を追加すれば
よい。

\begin{itemize}
\item 追加するコールの行使価格次第で，日経平均上昇時の正味損益関数の下落に
適当な\textbf{下限}を設定できる。追加したコール(行使価格 $K'$)より上では，
売ったコールの損失(傾き $-1$)と買ったコールの利益(傾き $+1$)が相殺し，
正味損益は水平になる。
\item ただし，コールのオプションプレミアムを\emph{支払う}分だけ，正味損益関数は
全体的に低下する。例えば行使価格68,000円のコールを追加した場合，この寄与は
1,200円程度である。
\item したがって，正味損益関数のピーク(レンジ内の一定収益)は，問題1のときより
約1,200円低くなる(約5,400円 $\to$ 約4,200円)。
\end{itemize}

正味損益関数の形は，$S_T<58{,}000$ で右上がり(下落側は元のまま)，レンジ内で
一定，$62{,}000<S_T<68{,}000$ で右下がり，$S_T>68{,}000$ で(マイナス水準で)
水平，となる。

\medskip
\textbf{備考。}解答例のグラフの詳細な数値は，Black-Scholesモデルを用いて，
無リスク金利 $r=3.0\%$，ボラティリティ $\sigma=30.0\%$ で計算されている。まだ
オプション価格の計算方法の話をしていないので，おおよその形があっていればよく，
オプション価格の分だけどの程度上下するかにはこだわらなくてよい。肝心なのは，
ポジションを追加することで，追加するオプションのペイオフが生じる以外の部分でも，
正味損益が(プレミアム支払の分だけ)全体的に低下することを理解することである。
\end{answerbox}
